\documentclass[12pt,a4paper]{article}
\usepackage[utf8]{inputenc}
\usepackage[T1]{fontenc}
\usepackage{amsmath, amssymb, amsthm}
\usepackage{hyperref}
\usepackage{geometry}
\geometry{margin=1in}

\title{Strict Resolution Architecture for the Erd\H{o}s Conjecture on Arithmetic Progressions of Primes}
\author{Charles EDOU NZE\thanks{Charles EDOU NZE, chercheur ind\'ependant}}
\date{\today}

\newtheorem{theorem}{Theorem}
\newtheorem{lemma}{Lemma}
\newtheorem{definition}{Definition}

\begin{document}
\maketitle

\begin{abstract}
This document lays out the foundational algebraic structures, rigorous algebraic proofs and structural formalism necessary to establish the existence of arbitrarily long arithmetic progressions of prime numbers.
\end{abstract}

\tableofcontents
\newpage

\section{Axiomatic Definitions and Type Specifications}

Let $\mathbb{P}$ denote the set of prime numbers. The conjecture, originally posed by Paul Erd\H{o}s and fully resolved by Ben Green and Terence Tao in 2004, states that the sequence of primes contains arithmetic progressions of arbitrary length.

\begin{definition}[Arithmetic Progression of Primes]
An arithmetic progression of length $k \in \mathbb{N}_{\geq 3}$ in $\mathbb{P}$ is a sequence of prime numbers $(p_1, p_2, \ldots, p_k)$ such that there exists a common difference $d \in \mathbb{N}$ where $p_i = p_1 + (i-1)d$ for all $1 \leq i \leq k$.
\end{definition}

\begin{definition}[Relative Upper Density]
Let $A \subset \mathbb{N}$. The upper density of $A$ relative to the primes $\mathbb{P}$ is defined as:
$$ \limsup_{N \to \infty} \frac{|A \cap [1, N]|}{|\mathbb{P} \cap [1, N]|} $$
\end{definition}

\section{Contextual Literature Research}

The cornerstone of this problem lies in additive combinatorics and analytic number theory:
\begin{itemize}
    \item \textbf{Szemer\'edi's Theorem (1975):} Any subset of the integers with positive upper density contains arithmetic progressions of arbitrary length.
    \item \textbf{The Green-Tao Theorem (2004):} Extended Szemer\'edi's theorem to sets with density zero, specifically the primes, by demonstrating that the primes form a pseudo-random subset of the integers.
\end{itemize}

An analogous breakthrough is the Gowers norms approach used in proving Szemer\'edi's theorem, which quantifies the pseudo-randomness of a set and controls the number of arithmetic progressions.

\section{Proof Strategy and Lemmas}

The proof relies on transferring Szemer\'edi's theorem to a sparse set (the primes) using a majorant measure.

\begin{lemma}[Transference Principle]
If $\nu: \mathbb{Z}_N \to \mathbb{R}_{\geq 0}$ is a pseudo-random measure and $f: \mathbb{Z}_N \to \mathbb{R}_{\geq 0}$ satisfies $0 \leq f \leq \nu$ and $\mathbb{E}(f) \geq \delta > 0$, then $f$ contains arithmetic progressions of length $k$.
\end{lemma}
\begin{proof}
Let $\nu$ be a measure satisfying the linear forms condition and the correlation condition (pseudo-randomness properties). We decompose the function $f$ as $f = f_{U} + f^{\perp}$, where $f_{U}$ is the anti-uniform (structured) component bounded by $1$, and $f^{\perp}$ is the uniform (pseudo-random) component with small Gowers $U^k$ norm.

Since $0 \leq f \leq \nu$, the generalized von Neumann theorem ensures that the contribution of $f^{\perp}$ to the count of arithmetic progressions of length $k$ is negligible. Thus, the count of progressions in $f$ is governed entirely by $f_{U}$.

Because $f_{U}$ is bounded ($0 \leq f_{U} \leq 1 + o(1)$) and retains the density $\mathbb{E}(f_{U}) \approx \mathbb{E}(f) \geq \delta$, Szemer\'edi's theorem directly applies to $f_{U}$. This implies $f_{U}$ contains $\gg \delta^{c_k} N^2$ arithmetic progressions. The negligible error from $f^{\perp}$ guarantees that $f$ itself contains arithmetic progressions of length $k$.
\end{proof}

\begin{lemma}[Construction of the Majorant]
There exists a pseudo-random measure $\nu$ that majorizes the primes, allowing the application of the Transference Principle.
\end{lemma}
\begin{proof}
Let $W = \prod_{p \leq w} p$ be the product of small primes up to $w$. We employ the $W$-trick to eliminate congruence obstructions. The modified prime counting function on a reduced residue class $b$ modulo $W$ is roughly proportional to the von Mangoldt function $\Lambda(Wn + b)$.

We construct the majorant $\nu(n)$ using Goldston-Pintz-Yildirim (GPY) truncated divisor sums:
$$ \nu(n) = \frac{\phi(W)}{W \log R} \left( \sum_{d | Wn + b, d \leq R} \mu(d) \log\left(\frac{R}{d}\right) \right)^2 $$
for an appropriate cutoff parameter $R = N^{\epsilon}$.

By Selberg sieve techniques, we verify that $0 \leq \Lambda'(n) \leq c \cdot \nu(n)$ for some constant $c>0$, where $\Lambda'(n)$ restricts to the primes. Furthermore, computing the moments of $\nu$ confirms it satisfies the linear forms and correlation conditions required for pseudo-randomness, completing the majorization.
\end{proof}

\section{Architecture for Autoformalization}

\begin{verbatim}
import Mathlib.Data.Nat.Prime
import Mathlib.Data.Finset.Basic
import Mathlib.Analysis.SpecialFunctions.Log.Basic

namespace ErdosArithmeticProgression

-- Definition of an arithmetic progression of length k in a set S
def HasArithmeticProgression (S : Set Nat) (k : Nat) : Prop :=
  Exists (fun a => Exists (fun d => d > 0 /\ \forall i, i < k → a + i * d \in S))

-- Formal statement of the Green-Tao Theorem
theorem green_tao_theorem (k : Nat) (hk : k \ge 3) :
  HasArithmeticProgression {p : Nat | p.Prime} k := by
  admit

-- The W-trick definition structure
def W_trick (w : Nat) : Nat :=
  -- Product of primes <= w
  admit

-- Definition of the pseudorandom majorant (GPY sieve)
def nu_majorant (n R W : Nat) : Real :=
  -- Truncated Moebius inversion squared
  admit

end ErdosArithmeticProgression
\end{verbatim}

\section*{References}
\begin{itemize}
    \item Green, B., \& Tao, T. (2008). \textit{The primes contain arbitrarily long arithmetic progressions}. Annals of Mathematics.
\end{itemize}

\end{document}
